\documentclass[10pt]{article}
\usepackage[margin=0.85in]{geometry}
\usepackage{amsmath,amssymb,amsthm,mathtools}
\usepackage{graphicx}
\usepackage{microtype}
\usepackage[hidelinks]{hyperref}

\newtheorem{theorem}{Theorem}
\newtheorem{remark}{Remark}
\newcommand{\ThetaSet}{\boldsymbol{\Theta}}
\newcommand{\N}{\mathbb N}

\title{Odd-$\kappa$ Symmetry of Finite-Order Threshold Sets}
\author{Candidate full answer to an open problem in arXiv:2201.09545}
\date{26 August 2026}

\begin{document}
\maketitle

\begin{center}
\fbox{\parbox{0.91\textwidth}{\textbf{Status.}
Candidate full solution, likely valid, pending human review.  For every
dimension $d\geq1$, every $\kappa\geq1$, and every $m\geq0$,
\[
 \ThetaSet_{m,\kappa}(\Delta)=-\ThetaSet_{m,\kappa}(\Delta).
\]
This supplies the odd-$\kappa$ case left open after Lemma 3.3 of the source.}}
\end{center}

\section{Source question}

Gol\'{e}nia and Mandich prove the displayed symmetry for even $\kappa$ and
ask whether it also holds for odd $\kappa$ \cite[p.~12, after Lemma 3.3]{GM}.

\begin{figure}[h]
\centering
\includegraphics[width=0.96\textwidth]{figures/open_problem_crop.png}
\caption{The open problem in the source, arXiv PDF page 12.}
\end{figure}

For $\vec x=(x_1,\ldots,x_d)\in[-1,1]^d$, the source defines
\begin{equation}\label{eq:g}
 g_{j\kappa}(\vec x)
   =\sum_{i=1}^d(1-x_i^2)U_{j\kappa-1}(x_i),
 \qquad j\in\N^*,
\end{equation}
where $U_n$ is the Chebyshev polynomial of the second kind.  The constant
energy surface is
\[
 S_E=\{\vec x\in[-1,1]^d:x_1+\cdots+x_d=E\}.
\]
For $m\geq1$, membership $E\in\ThetaSet_{m,\kappa}(\Delta)$ means that there
are $\vec x_0,\ldots,\vec x_m\in S_E$ and real weights
$\omega_0,\ldots,\omega_{m-1}\leq0$ such that
\begin{equation}\label{eq:witness}
 g_{j\kappa}(\vec x_m)
   =\sum_{q=0}^{m-1}\omega_q g_{j\kappa}(\vec x_q)
 \quad\text{for every }j\in\N^*.
\end{equation}
For $m=0$, the definition asks for $\vec x_0\in S_E$ with
$g_{j\kappa}(\vec x_0)=0$ for every $j$.

\section{Proof intuition}

Negate every coordinate of every witness point.  This changes the energy
from $E$ to $-E$.  The parity of $U_{j\kappa-1}$ gives a scalar sign which
may depend on $j$ when $\kappa$ is odd, but for each fixed $j$ it is identical
at every witness point.  Consequently it multiplies both sides of
\eqref{eq:witness} by the same number and cannot disturb the witness relation
or the signs of its weights.

\section{The theorem and proof}

\begin{theorem}\label{thm:main}
For all $d,\kappa\in\N^*$ and $m\in\N$,
\[
 \ThetaSet_{m,\kappa}(\Delta)=-\ThetaSet_{m,\kappa}(\Delta).
\]
In particular, Lemma 3.3 of \cite{GM} holds when $\kappa$ is odd.
\end{theorem}

\begin{proof}
The standard parity identity for Chebyshev polynomials is
\[
 U_n(-t)=(-1)^nU_n(t).
\]
Because $1-t^2$ is even, \eqref{eq:g} yields
\begin{equation}\label{eq:parity}
 g_{j\kappa}(-\vec x)
   =(-1)^{j\kappa-1}g_{j\kappa}(\vec x)
 \quad (j\in\N^*).
\end{equation}

Suppose first that $m\geq1$ and $E\in\ThetaSet_{m,\kappa}(\Delta)$, with
witnesses $\vec x_0,\ldots,\vec x_m$ and weights $\omega_q\leq0$ as in
\eqref{eq:witness}.  Put $\vec y_q=-\vec x_q$.  Then $\vec y_q\in S_{-E}$.
For every fixed $j$, multiply \eqref{eq:witness} by
$(-1)^{j\kappa-1}$ and use \eqref{eq:parity}; this gives
\[
 g_{j\kappa}(\vec y_m)
   =\sum_{q=0}^{m-1}\omega_qg_{j\kappa}(\vec y_q).
\]
The same weights are still nonpositive, so they witness
$-E\in\ThetaSet_{m,\kappa}(\Delta)$.

If $m=0$, choose $\vec x_0\in S_E$ with
$g_{j\kappa}(\vec x_0)=0$ for all $j$.  Equation \eqref{eq:parity} shows that
$-\vec x_0\in S_{-E}$ has the same zero property.  Thus in every case
$\ThetaSet_{m,\kappa}(\Delta)\subseteq-\ThetaSet_{m,\kappa}(\Delta)$.
Applying the same argument to $-E$ gives the reverse inclusion.
\end{proof}

\begin{remark}
For odd $\kappa$, the factor in \eqref{eq:parity} is $(-1)^{j-1}$ and hence
depends on $j$.  This causes no difficulty: the definition requires a
separate equality for each $j$, and within that equality the factor is common
to every witness point.  This is precisely the point that removes the parity
restriction in the source lemma.
\end{remark}

\section{Verification, scope, and novelty boundary}

The companion script \texttt{code/verify\_parity.py} constructs the $U_n$ by
their recurrence and symbolically verifies \eqref{eq:parity} for 224 triples
of dimension, $\kappa$, and $j$.  This is only a transcription diagnostic;
the proof above is exact and finite.

The theorem answers only the source's odd-$\kappa$ symmetry question for the
finite-order sets $\ThetaSet_{m,\kappa}(\Delta)$.  It does not identify the
full threshold set, prove its countability, or establish the conjectured
Mourre bands.

Bounded exact-sentence, arXiv-id, author, lemma-label, and keyword searches
through 26 August 2026 found the source and companion papers, but no later
primary source recording this answer or parity proof.  Priority is not
certified and expert review remains appropriate.

\section{Human review recommendation}

Check the source definition \eqref{eq:witness} and confirm that it is imposed
separately for each $j$.  Once this is verified, simultaneous negation and
Chebyshev parity prove the result without any analytic or convergence input.

\sloppy
\begin{thebibliography}{9}
\raggedright
\bibitem{GM}
S. Gol\'{e}nia and M.-A. Mandich,
\emph{Thresholds and more bands of a.c. Spectrum for the discrete
Schr\"odinger operator with a more general long range condition},
arXiv:2201.09545 (2022).
\end{thebibliography}

\end{document}
